%! TeX program = lualatex
% Copyright (c) 1996, 2004, 2006, 2009, 2014, 2016, 2019, 2023
% Tama Communications Corporation. All rights reserved.
%
% Redistribution and use in source and binary forms, with or without
% modification, are permitted provided that the following conditions
% are met:
% 1. Redistributions of source code must retain the above copyright
%    notice, this list of conditions and the following disclaimer.
% 2. Redistributions in binary form must reproduce the above copyright
%    notice, this list of conditions and the following disclaimer in the
%    documentation and/or other materials provided with the distribution.
%
% THIS SOFTWARE IS PROVIDED BY THE AUTHOR AND CONTRIBUTORS ``AS IS'' AND
% ANY EXPRESS OR IMPLIED WARRANTIES, INCLUDING, BUT NOT LIMITED TO, THE
% IMPLIED WARRANTIES OF MERCHANTABILITY AND FITNESS FOR A PARTICULAR PURPOSE
% ARE DISCLAIMED.  IN NO EVENT SHALL THE AUTHOR OR CONTRIBUTORS BE LIABLE
% FOR ANY DIRECT, INDIRECT, INCIDENTAL, SPECIAL, EXEMPLARY, OR CONSEQUENTIAL
% DAMAGES (INCLUDING, BUT NOT LIMITED TO, PROCUREMENT OF SUBSTITUTE GOODS
% OR SERVICES; LOSS OF USE, DATA, OR PROFITS; OR BUSINESS INTERRUPTION)
% HOWEVER CAUSED AND ON ANY THEORY OF LIABILITY, WHETHER IN CONTRACT, STRICT
% LIABILITY, OR TORT (INCLUDING NEGLIGENCE OR OTHERWISE) ARISING IN ANY WAY
% OUT OF THE USE OF THIS SOFTWARE, EVEN IF ADVISED OF THE POSSIBILITY OF
% SUCH DAMAGE.
%
%	rireki.tex	Version 3.1
%
%	URL: https://www.tamacom.com/rireki-j.html
%
%----------------------------------------------------------------------------
%
% 多摩通信社からのお知らせ
%
% 「履歴書猿人」(https://www.tamacom.com/engine/rireki2/) という履歴書作成のための
% Webアプリケーションを公開しております。TeX の知識なしにTeX の履歴書が作成でき、
% 様々なカスタマイズも可能です。こちらもご利用いただけると幸いです。
%
%----------------------------------------------------------------------------
%
% ヘッダー
%
\documentclass[b5j]{ltjsarticle}
\usepackage[deluxe,nfssonly]{luatexja-preset}
\usepackage{graphicx}
\usepackage{../../../../Styles/Rireki-Style}

\usepackage{datetime2}
\usepackage[
pdftitle={履歴書},
pdfauthor={アンドレ・モッシ},
]{hyperref}

% オプション
%
% 下記のオプションが利用可能です。

\空行挿入		% 学歴と職歴の間に空行を挿入します
%\性別欄なし		% 性別欄を削除します
%\写真欄なし		% 写真欄を削除します
\begin{document}

% ID情報

\姓{\LARGE Mossi Milla}
\名{\LARGE Roberto Andre}
\姓読み{ろばーと　あのどれ}
\名読み{もっし　みら}
\性別{Male}					% 男|女
\誕生日{2002年5月27日}
\現在日付{\the\year 年\the\month 月\the\day 日}
\年齢{\number\numexpr\the\year - 2002\relax}

% 顔写真
%
% 画像ファイルにはEPS フォーマット・縦横比4:3 のものをご使用ください。
% 縦を4cm に調整し、縦横比を変更せずに印刷します。
% 次のように指定します。
% \顔写真{photo.jpg}

\顔写真{../../../../Assets/2025PFPic.png}

% 現住所

\現住所郵便番号{16541}
\現住所{109 University Square, Erie, Pennsylvania, United States}
\現住所読み{ゆにばーしてぃ　すくえあ　えりー　ぺんしるばにあ}
\現住所市外局番{}
\現住所電話番号{}
\現住所呼び出し{}

% 連絡先

\連絡先郵便番号{}
\連絡先{\tt mossiroberto0392@gmail.com}
\連絡先読み{}
\連絡先市外局番{}
\連絡先電話番号{+1 814-790-1591}
\連絡先呼び出し{}

% 学歴、職歴
%
% 学歴、職歴を年月順に列挙してください。合計20個まで記入出来ます。
% 20個を超える部分は印刷されませんので、ご注意ください。
% 印刷順は、学歴=>職歴の順になります。

\学歴{2015}{4}{Dowal School　Enter School}      % {年}{月}{内容}
\学歴{2020}{3}{Dowal School　Graduated}
\学歴{2021}{8}{Gannon University Computer Science Admitted}
\学歴{2025}{5}{Gannon University Computer Science Graduated}
\学歴{2019}{4}{Appointed President of the Computer Club}
\学歴{2020}{3}{Won technical competition (Game Development category)}
\学歴{2025}{10}{Entered ISI Japanese Language School}

\職歴{2023}{6}{Joined Dinant Company as an intern}
\職歴{2023}{7}{Completed internship and left Dinant Company}

% 資格
%
% 資格を取得年月順に列挙してください。9つまで記入できます。
% 9つを超える部分は印刷されませんので、ご注意ください。

\資格{2023}{8}{Honduras Driver License}            % {取得年}{取得月}{資格}
\資格{2022}{12}{Harvard University\, CS50's Introduction to Computer Science}
\資格{2022}{5}{Extreme Networks\, Introduction to Future Networks}

% 個人情報

% 志望の動機と本人希望記入欄はlatex のコマンドを記述できます。

\志望の動機{
    My adventure started at 14 years old after viewing a video about game development in computer class.
    The video intrigued me, and I decided to start this journey into game development, not knowing the hardships
    I would have to go through. After winning my first competition showcasing my first-person shooter game made
    using C++ and Unreal Engine, I had found my passion for programming. From there on, I kept pushing forward
    every day, learning new languages, algorithms, and developing new projects. I won a second award at 18 years
    old with a survival game, developed web applications during my internship for stakeholders, and created a Neovim plugin
    for Copilot with more than 116 stars.

    Every new project I worked on increased my passion for programming. I hope to use all the skills I have
    acquired and my passion to help develop technologies that assist people, as well as take on new
    challenges that will contribute to your company's success.
}

\本人希望記入欄{
\begin{itemize}
	\item \textbf{Portfolio Website} https://andremossi.vercel.app/ \\
	      Portfolio webiste, were I talk about my achievements and show certificates

	\item \textbf{Linktree Website} https://andremossi-linktree.vercel.app/ \\
	      Linktree, were you can connect with me in different ways

	\item \textbf{Github} https://github.com/AndreM222 \\
	      Github projects

\end{itemize}
}

\サイン{Roberto Andre Mossi Milla}

\end{document}
